
En d'%
autres termes, \(K^{\bullet}(X)\) est la \(K^{\bullet}(S)\)-\(\lambda\)-algèbre engendrée par les générateurs \(F_1,\ldots,F_k\) soumis aux relations (cf. V 4.6)~:
\begin{enumerate}[label=\textup{\roman*)}]
\item \(\lambda^p(F_i)=0\) pour \(p>p_i\) ;
\item \(\sum_i F_i=E\).
\end{enumerate}

On commence par le cas où tous les \(p_i\) sont égaux à \(1\) ; le théorème s'énonce alors~:

\textbf{Corollaire 4.7.} \uline{Soient} \(S\) \uline{un schéma quasi}-\uline{compact}, \(\underline{E}\) \uline{un} \(\underline{O}_S\)\uline{-Module localement libre de rang} \(n\), \(X\) \uline{le fibré en drapeaux de} \(\underline{E}\), \(\xi_1,\ldots,\xi_n\) \uline{les classes des faisceaux inversibles} \(\underline{F}_1,\ldots,\underline{F}_n\) (\uline{notation de} 4.5). \uline{Alors} \(K^{\bullet}(X)\) \uline{est la} \(K^{\bullet}(S)\)-\uline{algèbre engendrée par les} \(\xi_i\), \uline{soumis aux relations}~:
\[
  \sigma_i(\xi_1,\ldots,\xi_n)=\lambda^i(E).
\]

On procède par récurrence sur \(n\). Pour \(n=1\), le théorème est trivial. Supposons qu'il soit montré pour \(n-1\) ; on considère alors \(\mathbf P(\underline{E})\), et sur \(\mathbf P(\underline{E})\) la suite exacte~:
\[
  0\longrightarrow \underline{F}\longrightarrow g^*(\underline{E})\longrightarrow \underline{O}_{\mathbf P(\underline{E})}(1)\longrightarrow 0,
\]
\(g\) étant le morphisme structural de \(\mathbf P(\underline{E})\). D'après le lemme 4.2, \(X\) s'identifie au fibré en drapeaux usuel \(\mathbf D(\underline{F})\) sur \(\mathbf P(\underline{E})\), les faisceaux canoniques se correspondant dans cette identification. Le diagramme~:
\[
\begin{tikzcd}[column sep=large,row sep=large]
\mathbf D(\underline{E})=\mathbf D(\underline{F}) \arrow[r,"h"] \arrow[dr,"f"'] & \mathbf P(\underline{E}) \arrow[d,"g"]\\
& S,
\end{tikzcd}
\]
montre que \(K^{\bullet}(S)\to K^{\bullet}(X)\) se factorise par \(K^{\bullet}(\mathbf P(\underline{E}))\). D'après l'hypothèse de récurrence appliquée à \(\mathbf D(\underline{F})\), et le corollaire 4.4, \(K^{\bullet}(X)\) admet pour base sur \(K^{\bullet}(\mathbf P(\underline{E}))\) les éléments de la forme~:
\[
  \xi_2^{\alpha_2}\ldots\xi_{n-1}^{\alpha_{n-1}},
  \qquad \text{avec } 0\le\alpha_i\le n-i;
\]
par ailleurs, \(K^{\bullet}(\mathbf P(\underline{E}))\) admet pour base sur \(K^{\bullet}(S)\) les éléments \(1,\xi_1,\ldots,\xi_1^{n-1}\) ; donc \(K^{\bullet}(X)\) admet pour base sur \(K^{\bullet}(S)\) les éléments~:
